\section{总结与延伸}

\subsection{核心知识回顾}

本讲系统介绍了深度生成建模的两大基础架构：

\textbf{1. Variational Autoencoder (VAE)}：
\begin{itemize}
    \item 通过编码器-解码器架构将数据压缩到概率性隐空间
    \item 损失函数 = 重建损失 + KL Divergence 正则化
    \item Reparameterization Trick 实现端到端可微训练
    \item 隐空间具有连续性和完整性，支持可解释的特征操控
\end{itemize}

\textbf{2. Generative Adversarial Network (GAN)}：
\begin{itemize}
    \item 生成器与判别器的 minimax 对抗博弈
    \item 不需要显式建模密度，直接学习分布变换
    \item 生成质量优异，但训练稳定性是核心挑战
    \item 衍生出 Progressive GAN、Conditional GAN、CycleGAN 等强大变体
\end{itemize}

\begin{importantbox}{生成模型的统一视角}
无论是 VAE 还是 GAN，它们的共同目标都是学习数据的底层概率分布。VAE 通过显式的密度估计实现这一目标，而 GAN 通过对抗训练隐式地学习数据分布。两种方法各有优劣，代表了生成建模的两种基本哲学：\textbf{显式建模} vs. \textbf{隐式采样}。
\end{importantbox}

\subsection{从 VAE/GAN 到现代生成模型}

本讲介绍的 VAE（2014）和 GAN（2014）是生成模型领域的两座里程碑。在它们之后，生成模型领域继续快速发展：

\begin{itemize}
    \item \textbf{Diffusion Models}（扩散模型，2020--）：通过逐步去噪过程生成数据，已成为图像生成的主流方法（如 Stable Diffusion、DALL-E 2/3）
    \item \textbf{Flow-based Models}（流模型）：通过可逆变换显式建模密度，如 Glow、RealNVP
    \item \textbf{Autoregressive Models}（自回归模型）：如 GPT 系列，逐 token 建模条件概率分布
    \item \textbf{Score-based Models}：通过估计数据分布的梯度场来生成样本
\end{itemize}

\begin{knowledgebox}{为什么 Diffusion Models 后来居上？}
Diffusion Models 结合了 VAE 和 GAN 的优点：它们像 VAE 一样有稳定的训练过程（基于似然优化），同时像 GAN 一样能生成高质量的样本。通过逐步去噪的马尔可夫过程，Diffusion Models 避免了 GAN 的训练不稳定性和 VAE 的生成模糊问题。MIT 6.S191 的后续讲座会专门讲解 Diffusion Models。
\end{knowledgebox}

\subsection{拓展阅读}

\begin{itemize}
    \item Kingma \& Welling, ``Auto-Encoding Variational Bayes,'' ICLR 2014 --- VAE 的开创性论文
    \item Goodfellow et al., ``Generative Adversarial Nets,'' NeurIPS 2014 --- GAN 的开创性论文
    \item Higgins et al., ``$\beta$-VAE: Learning Basic Visual Concepts with a Constrained Variational Framework,'' ICLR 2017
    \item Isola et al., ``Image-to-Image Translation with Conditional Adversarial Networks (Pix2Pix),'' CVPR 2017
    \item Karras et al., ``Progressive Growing of GANs,'' ICLR 2018
    \item Zhu et al., ``Unpaired Image-to-Image Translation using Cycle-Consistent Adversarial Networks (CycleGAN),'' ICCV 2017
    \item Brock et al., ``Large Scale GAN Training for High Fidelity Natural Image Synthesis (BigGAN),'' ICLR 2019
    \item Ho et al., ``Denoising Diffusion Probabilistic Models,'' NeurIPS 2020 --- Diffusion Models
    \item MIT 6.S191 课程主页：\href{https://introtodeeplearning.com}{introtodeeplearning.com}
    \item Lab 2 代码：\href{https://github.com/MITDeepLearning/introtodeeplearning}{github.com/MITDeepLearning/introtodeeplearning}
\end{itemize}

%% --- Fin del contenido --- %%

\end{document}
